\section{Methodology}\label{sec:method}

% Note: the stability-score walk-through and the calibration-sample rationale now live in Appendix D.

The strategy has two parts. I build an exposure score for every occupation from language-model ratings of its constituent tasks, then weight those scores by Scottish and rUK employment. The scores are the novel input, so most of this section concerns how they are constructed and how far they can be trusted; the labour-market accounting that follows is comparatively standard.

\subsection{Exposure scoring}\label{sec:scoring}

The OBR score each O*NET task from nought to a hundred for its amenability to AI over a ten-year horizon, separating tasks AI could substitute for from those it would merely complement. They classify an occupation as exposed where enough of its task content clears some threshold, and they map the US result onto UK occupations through a published crosswalk. I adopt that architecture, its scoring scale and its substitute--complement distinction, and depart from it in two ways. Their object is a single national figure, mine the difference between two regions. And they report point estimates without quantifying their dependence on the scoring framework, whereas I treat the scores as model-specific throughout and rest the regional claim on the component of the measurement that survives a change of model.
% check OBR threshold values (70, 40) 

\paragraph{Procedure.} I score the task content of every O*NET--SOC occupation $j$. For each task $t \in T_j$ the model returns two integer scores on $[0,100]$: a substitution score $s^{\mathrm{sub}}_{j,t}$, capturing the extent to which AI could perform the task without human involvement over the coming decade, and a complementarity score $s^{\mathrm{comp}}_{j,t}$, capturing the extent to which it could raise a worker's productivity in the task without replacing them. The prompt fixes the horizon, supplies the task description verbatim from O*NET, and directs the model to the considerations bearing on each judgement, such as whether the task requires physical presence, tacit empathy, or non-delegable accountability.

The two scores must then be combined into a single measure of how exposed a task is. For the exposure measurement I take the larger of the two, $s_{j,t} = \max(s^{\mathrm{sub}}_{j,t}, s^{\mathrm{comp}}_{j,t})$. I choose the maximum for commensurability, since it is the rule behind the OBR's single-amenability banding. It is not uniquely correct, and it is admittedly coarse. It treats a task scored $(90,10)$ as no less exposed than one scored $(90,90)$, and as the maximum of two noisy scores it carries a small upward bias that grows with the framing noise \citep{clark1961}. Two considerations mitigate this concern. First, the same aggregation rule is applied to each region, so by the logic of \cref{prop:cancellation} below its effect largely cancels in the regional differential. I verify this by reporting the gap under the substitution score alone and under a saturating alternative, $s_{j,t} = s^{\mathrm{sub}}_{j,t} + \bigl(1 - s^{\mathrm{sub}}_{j,t}/100\bigr)\, s^{\mathrm{comp}}_{j,t}$. This alternative treats complementarity as exposing only the share of a task left after substitution, and the resulting differentials remain closely aligned even where the underlying levels differ. Second, the productivity exercise in \cref{sec:specs} does not rely on the collapsed score. Instead, substitution and complementarity enter separately because they have opposite-signed effects in that specification.

Task scores then aggregate to the occupation. The occupation-level score is the importance-weighted mean of its task scores,
\begin{equation}\label{eq:aggregate}
E_j \;=\; \frac{\sum_{t \in T_j} \omega_{j,t}\, s_{j,t}}
               {\sum_{t \in T_j} \omega_{j,t}} \;\in\; [0,100],
\end{equation}
with O*NET importance weights $\omega_{j,t}$, so that an occupation's exposure reflects its central tasks rather than its peripheral ones. Because the scores attach to US occupations, I map $E_j$ onto a UK SOC 2020 score $E^{\mathrm{UK}}_k$ for minor group $k$ through an established crosswalk, working at the three-digit minor-group level so that Scottish employment cells stay populated when survey waves are pooled.

Classification converts the continuous scores into the exposed and unexposed categories. My baseline thresholds are those of the OBR, $(\theta^{\mathrm{sub}}, \theta^{\mathrm{comp}}) = (70, 40)$, applied at the task level: a task is substituted where $s^{\mathrm{sub}}_{j,t} \ge \theta^{\mathrm{sub}}$, complemented where $s^{\mathrm{comp}}_{j,t} \ge \theta^{\mathrm{comp}}$ and the task is not substituted, and unexposed otherwise. The importance-weighted shares of task content in the two exposed classes, $\mathrm{Sub}_j$ and $\mathrm{Comp}_j$, then classify the occupation. Occupation $j$ is materially exposed where $\mathrm{Sub}_j + \mathrm{Comp}_j \ge \tfrac{1}{2}$, and an exposed occupation is substituted or complemented according to the majority of its exposed mass. Where my rule departs is that the OBR band a single amenability score, whereas I elicit the two channels separately and let each face its own threshold. The Scottish composition question turns on the substitution--complementarity margin, and eliciting the two channels separately measures that margin directly, where bands of a single amenability score can only imply it. Sensitivity of the classified indices to the threshold pair is reported at $(60,30)$ and $(80,50)$ alongside the baseline.

The resulting score is common to both regions. O*NET records the tasks of an occupation, not of an occupation in a place, so the same $E^{\mathrm{UK}}_k$ applies in Scotland and the rUK. To guard against the prompt-sensitivity this literature has repeatedly flagged \citep{eloundou2024, yin2026}, I fix and freeze a single model and prompt before the main run, so that the only variation across occupations is genuine.\footnote{I record a SHA-256 hash of the frozen system message and user-message template together with every score, so any later change to either is detectable and the figures reproduce from the frozen configuration.}

\subsection{Regional indices and the compositional structure of the gap}\label{sec:indices}

The question the Scottish Fiscal Commission needs answered is whether Scotland's workforce sits more or less in the path of this technology than the rest of the UK's. The natural summary is an employment-weighted exposure index for each region: how exposed is the average job, given where each region's employment actually sits across occupations. For region $r \in \{\Scot, \rUK\}$, let $\eta_{k,r} = l^{\mathrm{pool}}_{k,r} \big/ \sum_{k'} l^{\mathrm{pool}}_{k',r}$ be the pooled employment share of minor group $k$. I form a classified exposed-employment share and a continuous index,
\begin{equation}\label{eq:indices}
\Ehat_r = \sum_k \eta_{k,r}\,
  \ind{\mathrm{Class}(k) \in \{\mathrm{Sub}, \mathrm{Comp}\}},
\qquad
\Ebar_r = \sum_k \eta_{k,r}\, \frac{E^{\mathrm{UK}}_k}{100}
  \;\in\; [0,1].
\end{equation}
The first answers what fraction of employment sits in materially exposed occupations. The second uses the full scores and is the object for which the formal results below are exact. The gap $\Delta\Ehat = \Ehat_{\Scot} - \Ehat_{\rUK}$, with continuous analogue $\dEbar$, is the object of interest.

A shift-share decomposition makes precise what this gap can and cannot contain. Any regional difference in an employment-weighted index has, in principle, two sources: the regions may distribute their employment differently across occupations, or the same occupation may be differently exposed in the two places. Writing the gap as
\begin{equation}\label{eq:shiftshare}
\Delta\Ehat =
\underbrace{\sum_k \bigl(\eta_{k,\Scot} - \eta_{k,\rUK}\bigr)\,
  \ind{\mathrm{Exp}(k)}_{\rUK}}_{\text{composition}}
\;+\;
\underbrace{\sum_k \eta_{k,\Scot}\,
  \bigl(\ind{\mathrm{Exp}(k)}_{\Scot} -
        \ind{\mathrm{Exp}(k)}_{\rUK}\bigr)}_{\text{within-occupation}}
\end{equation}
separates the two sources. Since the same score applies on both sides of the border, the within-occupation term vanishes identically. The entire gap is therefore compositional. Any Scotland--rUK differential reflects Scotland's distinctive occupational mix, with more public administration and caring work and less professional and financial services.

This compositional reading rests on one assumption: the within-occupation task profile of each UK minor group is the same in both regions, so that a single $E^{\mathrm{UK}}_k$ describes minor group $k$ on either side of the border. The assumption is not that occupations are internally homogeneous; \citet{henseke2025} show they are not, with the highest score dispersion in cells where a generic task description admits both an exposed and an unexposed reading. It is the weaker claim that whatever heterogeneity exists is distributed similarly in the two regions.

\begin{assumption}[Region invariance]\label{as:reginv}
Writing $T$ for a worker's task bundle, $k$ for their minor group and $R \in \{\Scot, \rUK\}$ for region, the conditional distribution of tasks is region-invariant,
\begin{equation}\label{eq:reginv}
\Pr(T \in A \mid k, R = \Scot) \;=\; \Pr(T \in A \mid k, R = \rUK)
\quad \text{for all } A .
\end{equation}
\end{assumption}

\noindent
Internal homogeneity would require $T$ to be degenerate given $k$. What I need is only \cref{eq:reginv}, so that whatever within-occupation heterogeneity exists is distributed identically on either side of the border. Even this overstates what the differential requires: writing $e(\cdot)$ for task-level exposure, only the conditional means need agree, $\E[e(T) \mid k, \Scot] = \E[e(T) \mid k, \rUK]$, and any common departure of that mean from the O*NET-measured profile cancels from $\Delta\Ehat$ in any case. Within one national economy this is reasonable. The qualifications, technologies, and regulatory standards that define an occupation are set at UK level, and the coarse working level of the crosswalk averages over the fine distinctions along which any residual regional difference would appear \citep{autor2013}. The assumption fails only if task content differs across regions in a way correlated with exposure. I cannot test this directly without region-specific task data, but the international evidence suggests such within-occupation differences are second-order relative to composition \citep{handel2012, autor2013}. 
% And where the assumption matters, it is needed only in differenced form: any error in $E^{\mathrm{UK}}_k$ common to both regions cancels from $\Delta\Ehat$, exactly as the model bias treated next does.

Finally, a single $E^{\mathrm{UK}}_k$ attaches to a minor group irrespective of the industry in which the occupation is employed. As a result, any occupation-by-industry variation in exposure, between the same accountant working in a bank and in a local authority, say, is removed from the regional differential by construction. The design can still identify whether the Scotland--rUK gap arises from differences in industry composition or from differences in occupational composition within industries, and \cref{sec:anatomy} decomposes the differential along those lines. What it cannot capture is variation in exposure within a given occupation--industry cell across regions. Doing so would require industry-specific task data, which no available source provides.

\subsection{Scoring error and the robustness of the differential}\label{sec:noise}

A task score can be wrong for more than one reason, and the design responds to each of them separately. Meaning-preserving changes to the prompt move the score, and that framing noise is what the calibration of \cref{eq:variance} measures. Each model carries a persistent tendency to over- or under-score, and that systematic bias \citep{yin2026} is what \cref{prop:cancellation} addresses and the cross-model test of \cref{sec:crossmodel} probes. Aggregation and crosswalk error is bounded by the operator and threshold robustness, and sampling error in the employment weights by the APS bootstrap; both are handled where they arise.

The cross-model divergence documented by \citet{yin2026} would be fatal to a study reporting absolute levels as objective measurements. It is far less threatening to a study of regional differences, however, and this subsection formalises why. Consider a score generated by model-$M$, decomposed as
\begin{equation}\label{eq:decomposition}
s_j \;=\; \theta^{\star}_j + b^{M}_j + \varepsilon_j,
\qquad
\varepsilon_j \sim N\bigl(0, \sigma^2_{g(j)}\bigr),
\end{equation}
where $\theta^{\star}_j$ is the unobserved cross-model consensus, $b^{M}_j$ is model $M$'s systematic bias, meaning its persistent tendency to over or under-score occupation $j$, and $\varepsilon_j$ is mean-zero framing noise indexed by the occupational family $g(j)$ defined in the calibration below. The consensus $\theta^{\star}_j$ is a device for organising the decomposition, not an estimand: only contrasts between models are ever observed, and nothing below requires access to it. The two error terms play different roles and therefore require different treatment. The noise term $\varepsilon_j$ is the component that averaging and simulation can reduce. $b^{M}_j$, by contrast, is the persistent model-specific bias identified by \citet{yin2026}, affecting every absolute estimate a model produces. I measure the former directly and show that the latter cancels in the regional differential.

Conditioning on the frozen model absorbs $b^{M}_j$ into a model-specific target $\theta^{M}_j = \theta^{\star}_j + b^{M}_j$, leaving the classical measurement model $s_j = \theta^{M}_j + \varepsilon_j$, with group variances $\sigma^2_{g(j)}$ to be calibrated below. Writing the continuous UK score as $E^{M}_k = E^{\star}_k + b^{M}_k$, the cancellation result follows.

\begin{proposition}[Common-mode cancellation]\label{prop:cancellation}
The model-induced error in the differential $\dEbar^{M} = \Ebar^{M}_{\Scot} - \Ebar^{M}_{\rUK}$ is
\begin{equation}\label{eq:prop1}
\dEbar^{M} - \dEbar^{\star} \;=\;
\frac{1}{100} \sum_k \bigl(\eta_{k,\Scot} - \eta_{k,\rUK}\bigr)\,
b^{M}_k .
\end{equation}
Hence (i) a pure level shift $b^{M}_k = b^{M}$ leaves the differential exact, $\dEbar^{M} = \dEbar^{\star}$; and (ii) in general only the part of the bias co-varying with the regional employment-share gap survives, with the mean bias being removed by construction.
\end{proposition}

\noindent
\Cref{app:proofs} proves \cref{prop:cancellation} and states the demeaning step it turns on.

The absolute index $\Ebar^{M}_r$ carries contamination of the order of the mean bias, which is what the three-fold divergence of \citet{yin2026} measures. The differential retains only the part of that bias co-varying with the cross-region employment-share gap. So the same compositional weighting that zeroes the within-occupation term of \cref{eq:shiftshare} also removes the common-mode bias, and the differential is robust where the levels are fragile. How much of the bias in fact survives is an empirical matter, and \cref{sec:crossmodel} measures it.

The proposition also says something a re-scoring experiment can check, and I record that implication as a corollary so the cross-model test of \cref{sec:crossmodel} has a formal target to be judged against.

\begin{corollary}[Cross-model stability]\label{cor:crossmodel}
For any two models $M$ and $M'$,
\begin{equation}\label{eq:cor1}
\dEbar^{M} - \dEbar^{M'} \;=\;
\frac{1}{100} \sum_k \Delta\eta_k
\bigl(b^{M}_k - b^{M'}_k\bigr),
\qquad
\bigl|\dEbar^{M} - \dEbar^{M'}\bigr| \;\le\;
\frac{1}{100}\,\bigl\lVert \Delta\eta \bigr\rVert_2\,
\bigl\lVert b^{M} - b^{M'} \bigr\rVert_2 ,
\end{equation}
where $\Delta\eta_k = \eta_{k,\Scot} - \eta_{k,\rUK}$ and the bias differences are demeaned. The cross-model movement in the differential is zero under a common level shift and bounded by the alignment of the bias difference with the share gap in general.
\end{corollary}

\noindent
\Cref{eq:decomposition} is additive, so \cref{prop:cancellation} protects the differential against an additive bias only. A model may instead disagree about the \emph{scale} of the index, compressing or stretching the score distribution while preserving the ordering. The next corollary covers that case, which \cref{sec:crossmodel} shows is the one the data present.

\begin{corollary}[Affine bias and scale-free differentials]\label{cor:affine}
Suppose model $M$'s scores are affine in the consensus up to an idiosyncratic term,
$E^{M}_k = a^{M} + c^{M} E^{\star}_k + b^{M}_k$ with $c^{M} > 0$. Then
\begin{equation}\label{eq:cor2}
\dEbar^{M} \;=\; c^{M}\, \dEbar^{\star} \;+\;
\frac{1}{100} \sum_k \Delta\eta_k\, b^{M}_k .
\end{equation}
Hence (i) the intercept $a^{M}$ cancels exactly, by \cref{eq:A4}, and the slope rescales the differential proportionally rather than cancelling; (ii) writing $\mathrm{SD}^{M}$ for the employment-weighted standard deviation of $\{E^{M}_k\}$, the standardised differential $\dEbar^{M}/\mathrm{SD}^{M}$ removes both $a^{M}$ and $c^{M}$, exactly when $b^{M}_k \equiv 0$ and up to the residual term otherwise; and (iii) replacing each score by its employment-weighted percentile rank yields a differential invariant to \emph{any} strictly increasing transformation of the scale, of which the affine map is a special case.
\end{corollary}

\noindent
A statement about the differential is \emph{scale-free} when it does not depend on how widely the index happens to be spread under a given model. Such a statement locates Scotland relative to the dispersion of UK occupations, and not at a distance in one model's units. \Cref{cor:affine} offers two such statements because they buy different things. Dividing by the model's own dispersion keeps the differential cardinal, so it remains comparable with other standardised gaps, but undoes only an affine distortion. Ranks undo any monotone distortion whatever, at the price of discarding everything but order.

The results answer different failures. \Cref{prop:cancellation} handles a model that reads the whole economy as more or less exposed, \cref{cor:affine} one that uses more or less of the available scale, and the percentile form one that agrees only about the ordering, the last being non-parametric and correspondingly the weakest statement about magnitude. That is, re-drawing the scores from a different model should move the levels, rescale the differential, and leave its sign and scale-free magnitude where they stood. I test exactly this by re-scoring the full task set on two alternative models and comparing the cross-model stability of $\dEbar$ against that of $\Ehat_{\Scot}$ and $\Ehat_{\rUK}$ individually.\footnote{I use one model from a different provider (OpenAI's GPT-4o) and one from a different tier of the same provider (Anthropic's Claude Haiku 4.5).} If the levels move while the gap holds, the regional finding survives the critique that undermines the absolute numbers.

\paragraph{Calibration and propagation.} The differencing handles the bias but the noise has to be measured. Framing noise is the variation in a score under changes to the prompt that leave its meaning intact. I do not assume a functional form for it. I estimate its variance from a stratified experiment. The experiment draws $N=90$ occupations, eighteen from each of the five broad occupational families obtained by collapsing the O*NET major groups onto their leading digit: ten selected at random to estimate the variance and eight nearest the classification boundary, held out for the stability check. The five-way grouping is coarse deliberately. It is the finest partition at which every cell carries enough occupations for a variance estimated from four scoring conditions to be usable, and $g(j)$ denotes it throughout.
%\footnote{Estimating the variance on a purely boundary-proximate sample would bias it upward if ambiguity rises towards the middle of the scale, which is where the boundaries sit; the random stratum avoids this, and the boundary stratum is used only for out-of-sample validation.} 
Each sampled occupation's full task set is re-scored under the conditions of \cref{tab:calibration}. Four of these conditions (the baseline, the assessment criteria in reverse order, a paraphrased rubric of identical content, and a reordered task presentation) are meaning-preserving perturbations of the frozen prompt, and the disagreement among them is what identifies the noise. Writing $K_w$ for this set, with $|K_w| = 4$,
\begin{equation}\label{eq:variance}
\widehat{\sigma}^2_j \;=\;
\frac{1}{|K_w| - 1} \sum_{k \in K_w}
\bigl(s^{(k)}_j - \bar{s}_j\bigr)^2,
\qquad
\widehat{\sigma}^2_g \;=\;
\frac{1}{|J_g|} \sum_{j \in J_g} \widehat{\sigma}^2_j ,
\end{equation}
where $J_g$ collects the random-stratum occupations of family $g$. Two further conditions, a five-year horizon and a single composite score, change the question rather than its phrasing, and a shorter horizon is not an error in measuring the ten-year one. I therefore report their occupation-level shifts as named sensitivities and do not pool them into $\widehat{\sigma}^2_g$.\footnote{Pooling them would conflate framing noise with a change of estimand and overstate the variance to which the propagation applies. By \cref{prop:cancellation}, the component of either shift common across occupations is in any case innocuous for the differential.}

\begin{table}[t]
\centering
\caption{Calibration design: scoring conditions and the role of each}
\label{tab:calibration}
\small
\begin{tabular}{lll}
\toprule
Variant & What changes & Role \\
\midrule
V0 & Baseline prompt, 10-year horizon & noise (reference) \\
V1 & Assessment criteria in reverse order & noise: order and anchoring \\
V2 & Paraphrased rubric, identical content & noise: wording \\
V3 & Task presented before occupation & noise: framing \\
V4 & 5-year horizon, else identical & estimand sensitivity: horizon \\
V5 & Single composite score & estimand sensitivity: rubric form \\
\bottomrule
\end{tabular}
\begin{minipage}{0.92\textwidth}
\vspace{0.4em}\footnotesize
The within-model variance of \cref{eq:variance} is estimated from V0, V1, V2 and V3 on the random sample only.; V4 and V5 are reported as named sensitivities. Cross-model contrasts come from the full re-scores under the two alternative models, not from a calibration condition.
\end{minipage}
\end{table}

With the variance in hand, I propagate it into the continuous aggregates by simulation. Across $R = 1{,}000$ draws I perturb each occupation score, $s^{(r)}_j \sim N\bigl(s_j, \widehat{\sigma}^2_{g(j)}\bigr)$ truncated to $[0,100]$, re-aggregate to UK minor groups, re-weight, and record $\dEbar^{(r)}$. I report the resulting 95\% interval and the share of draws in which the gap keeps its sign.\footnote{A gap whose interval excludes zero is robust to scoring noise, and one whose interval includes zero is not. As a check I allow intra-group correlated errors, $\varepsilon^{(r)}_g \sim N\bigl(0, \widehat{\sigma}^2_g [(1-\rho)I + \rho \bm{1}\bm{1}^{\top}]\bigr)$ for $\rho \in \{0, 0.2, 0.4\}$, where $\rho$ is the degree to which scoring errors are assumed to be correlated within an occupational family, so the higher values are the cautious cases. Correlated framing errors would otherwise lead independent draws to understate aggregate uncertainty. \cref{app:proofs} derives the resulting inflation factor.} The uncertainty propagation is carried out on the continuous index by design. It is the object for which \cref{prop:cancellation} holds exactly and for which the noise model in \cref{eq:variance} is defined. Alongside the scoring draws, I bootstrap the pooled employment shares over persons within region and year and combine the two independent sources of uncertainty to construct an interval for $\dEbar$. To show their relative importance, I also report the scoring and sampling intervals separately.

The classified share needs a separate treatment, because it is nonlinear at the thresholds. A small perturbation to a task sitting near a cut-off can flip its class, and with it the classification of an occupation. I therefore give each occupation a stability score: the importance-weighted probability that its task content keeps its class under re-framing, obtained by reading each task's distance from the nearest cut-off in units of the calibrated noise standard deviation and passing that distance through the normal cdf. \Cref{app:calibration} states the construction and the reporting landmark I flag against. 
In the event the flag binds nowhere. No occupation and no minor group falls below the landmark, so \cref{sec:fragility} compares the held-out boundary stratum against the random one instead of comparing flagged occupations against unflagged ones, and the stability apparatus enters as a null diagnostic.

\subsection{Empirical specifications}\label{sec:specs}

With exposure scores and regional weights in hand, I put the measured differential to work on three objects the Scottish Fiscal Commission forecasts, carrying the scoring uncertainty of \cref{sec:noise} through into each. The productivity projection is the one strictly necessary to the fiscal question, since it alone converts an exposure gap into the units in which a block-grant-adjustment claim would have to be made. The employment gradient restates the compositional reading of \cref{eq:shiftshare} in observable quantities, and the wage gradient is the micro counterpart of the aggregate propagation, where the calibrated scoring variance meets a regressor measured with error. None is offered as a causal estimate.

The first step expresses the compositional gap as an estimable gradient, regressing log occupational employment, and in a second specification the employment share, on the continuous exposure score, a Scotland indicator and their interaction, within one-digit SOC major groups and weighted by cell employment. The interaction asks whether, within broad groups, Scottish employment loads differentially onto exposed work: a negative estimate would mean composition insulates Scotland, a positive one the reverse. \Cref{app:specifications} gives the specification in full and \cref{sec:anatomy} reports the estimates, which restate the decomposition of \cref{eq:shiftshare} in observable quantities and add nothing to it.

The second step translates the exposure differential into productivity, using the task-based framework of \citet{acemoglu2025} as parameterised by the \citet{obr2025}. In that framework, AI raises total factor productivity by lowering the cost of the tasks it can perform, so the gain to a region depends on how much of its employment sits in exposed work. The proportional TFP gain for region $r$ is
\begin{equation}\label{eq:tfp}
\Delta \ln \mathrm{TFP}_r \;\approx\;
\kappa \sum_k \eta_{k,r}\, \pi_k\, \Ebar_k ,
\end{equation}
with $\pi_k$ the task-cost share and $\kappa$ the calibrated average cost saving on exposed tasks. Substitution and complementarity carry different economics here. The first is a cost saving on tasks the technology takes over; the second augments the worker who is retained. I therefore run the projection on the two channels separately rather than on their maximum, mapping the substitution channel to the cost-saving term above and reporting the complementarity channel alongside it. I also report the composite version for comparability with the OBR exercise. I import $\kappa$ and the horizon parameters from the OBR/Acemoglu calibration rather than estimating them, so the Scottish and rUK paths differ only through their employment-weighted exposure. \citet{acemoglu2025} builds $\kappa$ from the two task-level randomised experiments of \citet{noy2023} and \citet{brynjolfsson2025}, averaging their labour cost savings of 40\% and 14\% to about 27\%, then converting to overall cost savings using industry labour shares; the baseline calibration sets the task-cost weight to $\pi_k = \pi = 0.23$, his share of exposed tasks judged profitably automatable over the decade, uniform across occupations. Because $\pi_k$ is constant in $k$, the regional differential in each channel's implied TFP path inherits the cancellation of \cref{prop:cancellation} exactly, by \cref{cor:functionals}; only the levels of the path carry the model dependence. \Cref{app:specifications} derives \cref{eq:tfp} and gives the calibration in full, and I report the implied output-per-worker gap as a band reflecting the OBR's own parameter range.\footnote{The contribution is the regional differential in the implied path, not a reassessment of Acemoglu's headline magnitude, which is contested and outside the scope of this dissertation.}

The third step estimates the exposure--wage relationship at worker level,
\begin{equation}\label{eq:wage}
\ln w_{i,t} = \alpha + \beta_1 E_{j(i)} + \beta_2\, \Scot_i +
\beta_3 \bigl( E_{j(i)} \times \Scot_i \bigr) + \gamma X_{i,t} +
\delta_s + \delta_t + \varepsilon_{i,t},
\end{equation}
where $\ln w_{i,t}$ is the log gross hourly pay of worker $i$ in survey year $t$, $E_{j(i)} \in [0,1]$ is the continuous exposure score of the worker's minor group $j(i)$, common to both regions and constant within an occupation, and $\Scot_i$ indicates residence in Scotland. The coefficient $\beta_1$ is the rUK gradient of pay in exposure, $\beta_2$ the Scotland wage-level shift, and $\beta_3$, the object of interest, the Scotland-specific rotation of that gradient. The controls $X_{i,t}$ are age, its square, sex and full-time status; $\delta_s$ and $\delta_t$ are industry-section and survey-year fixed effects; and observations are weighted by the APS income weight. Two further controls enter the richer specifications of \cref{tab:wagespecs}. The first is highest qualification, the classic omission of the earnings-function literature: exposed occupations are graduate-heavy, so any Scotland--rUK difference in the within-occupation qualification profile would otherwise load onto $\beta_3$. The second is a public-sector indicator, because I read the estimated rotation partly through the public-sector pay premium \citep{cribb2025} and Scotland's public sector is proportionally larger, so I test that reading against the indicator. $E_{j(i)}$ varies only across occupations, every worker in a minor group shares one value of the regressor and the within-group correlation of residuals inflates conventional standard errors by the Moulton factor \citep{moulton1990}. I therefore cluster by occupation, the level at which the regressor varies.